% Vietnamese translation for OI Public Library
% Source-File: IOI中国国家候选队论文/2006/冯威/冯威.ppt
\documentclass[11pt]{article}
\usepackage{oipl-vn}

\newcommand{\Oh}{\mathcal{O}}

\title{Hệ ràng buộc sai phân\\{\large Ghi chú theo slide}}
\author{Feng Wei}
\date{2006}

\begin{document}
\maketitle

\origtitle{Difference constraints}
\sourcepath{IOI China candidate team papers / 2006 / Feng Wei / Feng Wei.ppt}

\section{Phạm vi}

Bộ trình chiếu đi kèm bài viết về hệ ràng buộc sai phân. Ghi chú này dựa trên
bản bài viết đã dịch, trong đó trọng tâm là chuyển một hệ bất đẳng thức giữa
hiệu hai biến thành bài toán đường đi trên đồ thị.

\section{Từ bất đẳng thức sang cạnh}

Một ràng buộc dạng
\[
  x_u-x_v\ge w
\]
có thể xem như cạnh \(v\to u\) trọng số \(w\). Khi tìm nhãn đỉnh lớn nhất phù
hợp với các cạnh, thao tác nới lỏng cạnh yêu cầu
\[
  d[u]\ge d[v]+w.
\]
Vì thế một hệ nhiều bất đẳng thức hiệu hai biến được chuyển thành một đồ thị
có trọng số. Nếu quá trình nới lỏng chỉ ra một chu trình làm giá trị phải tăng
mãi, hệ tương ứng không có nghiệm hữu hạn theo mục tiêu đang xét.

\section{Vai trò của Bellman--Ford}

Bellman--Ford là công cụ phù hợp vì nó không đòi trọng số không âm và có cơ
chế phát hiện chu trình bất thường. Trong thực hành, nhiều bài dùng biến thể
hàng đợi để tránh duyệt lại mọi cạnh ở mỗi lượt, nhưng mô hình đúng vẫn là
điểm quan trọng nhất.

\section{Ví dụ điển hình}

Với bài chọn số trên các đoạn, đặt tổng tiền tố \(S_i\). Điều kiện trong đoạn
\([a,b]\) có ít nhất \(c\) số được chọn trở thành
\[
  S_b-S_{a-1}\ge c.
\]
Điều kiện mỗi vị trí được chọn hoặc không được chọn cho ta
\[
  0\le S_i-S_{i-1}\le 1.
\]
Tất cả đều là ràng buộc sai phân, nên có thể đưa vào cùng một đồ thị.

\section{Thông điệp}

Giá trị của phương pháp nằm ở việc nhận ra biến tiền tố hoặc biến tích lũy.
Một khi các điều kiện được viết thành hiệu của hai biến, phần thuật toán còn
lại trở thành một bài đường đi chuẩn.

\end{document}
